\subsection{Convex functions}

DEFINITION 5. --- Let X be a convex subset of the affine space E. A real-valued finite function, defined over X is convex (resp. strictly convex) if for any two distinct points x, y of X and any real number \(\lambda\) , \(0 < \lambda < 1\) , we have :

\[
f (\lambda x + (1 - \lambda) y) \leqslant \lambda f (x) + (1 - \lambda) f (y)\tag{1}
\]

(resp.

\[
f (\lambda x + (1 - \lambda) y) <   \lambda f (x) + (1 - \lambda) f (y)).\tag{2}
\]

When E = R, this definition of convex function is the same as that in FVR, I, p.~32. Further, f is convex (resp. strictly convex) in X if, and only if, for every affine line \(D \subset E\), the restriction of f to \(X \cap D\) is convex (resp. strictly convex) in \(X \cap D\).

Examples. --- If \(f\) is an affine linear function on \(\mathbf{E}\) , then \(f\) and \(f^2\) are convex functions on \(\mathbf{E}\) ; this is obvious for \(f\) since

\[
f (\lambda x + (1 - \lambda) y) = \lambda f (x) + (1 - \lambda) f (y);
\]

on the other hand, if \(\alpha = f(x)\) , \(\beta = f(y)\) , then;

\[
\lambda \alpha^ {2} + (1 - \lambda) \beta^ {2} - (\lambda \alpha + (1 - \lambda) \beta) ^ {2} = \lambda (1 - \lambda) (\alpha - \beta) ^ {2} \geqslant 0
\]

for \(0 < \lambda < 1\) ; further, the restriction of \(f^2\) to an affine line \(\mathbf{D} \subset \mathbf{E}\) is strictly convex if \(f|\mathbf{D}\) is not a constant.

A real-valued function \(f\) , defined over \(\mathbf{X}\) , is concave (resp. strictly concave) if \(-f\) is convex (resp. strictly convex). That is to say, for every two distinct points \(x, y\) of \(\mathbf{X}\) and every number \(\lambda\) , such that \(0 < \lambda < 1\) , we have

\[
f (\lambda x + (1 - \lambda) y) \geqslant \lambda f (x) + (1 - \lambda) f (y)
\]

(resp.

\[
f (\lambda x + (1 - \lambda) y) > \lambda f (x) + (1 - \lambda) f (y)).
\]

A mapping of X in R is affine if it is both convex and concave (cf.~II, p.~78, exerc. 11).

PROPOSITION 19. --- Let X be a convex set of the affine space E; and let f be a real-valued function defined over X. Denote the set of points \((x, a) \in \mathbf{E} \times \mathbf{R}\) for which \(x \in X\) and \(f(x) \leqslant a\) (resp. \(x \in X\) and \(f(x) < a\) ) by F (resp. F'). Then the following conditions are equivalent:

\begin{enumerate}
\def\labelenumi{\alph{enumi})}
\item
  The function \(f\) is convex.
\item
  The set \(\mathbf{F}\) in the affine space \(\mathbf{E} \times \mathbf{R}\) is convex.
\item
  The set \(\mathbf{F}'\) in the affine space \(\mathbf{E} \times \mathbf{R}\) is convex.
\end{enumerate}

We show that \(a) \Rightarrow c)\) . Let \((x, a)\) and \((y, b)\) be two points of \(F'\) and \(0 < \lambda < 1\) , then \(f(x) < a\) , \(f(y) < b\) and if \(f\) is convex

\[
f (\lambda x + (1 - \lambda) y) \leqslant \lambda f (x) + (1 - \lambda) f (y) <   \lambda a + (1 - \lambda) b
\]

which shows that the point \(\lambda(x, a) + (1 - \lambda)(y, b)\) of \(\mathbf{E} \times \mathbf{R}\) belongs to \(\mathbf{F}'\) . Thus \(\mathbf{F}'\) is convex.

Next we show that \(c) \Rightarrow b)\) . If \((x, a), (y, b)\) are two points of \(F\) then for every \(\varepsilon > 0\) , \((x, a + \varepsilon)\) and \((y, b + \varepsilon)\) belong to \(F'\) and, if \(0 < \lambda < 1\) , the same is true of \((\lambda x + (1 - \lambda)y, \lambda a + (1 - \lambda)b + \varepsilon)\) ; by the definition of \(F\) this implies that \((\lambda x + (1 - \lambda)y, \lambda a + (1 - \lambda)b)\) belongs to \(F\) .

Finally \(b) \Rightarrow a)\) , for (with the above notation), if \((\lambda x + (1 - \lambda)y, \lambda a + (1 - \lambda)b)\) belongs to F then

\[
f (\lambda x + (1 - \lambda) y) \leqslant \lambda a + (1 - \lambda) b
\]

provided \(a \geqslant f(x)\) and \(b \geqslant f(y)\) ; hence (1) follows and \(f\) is convex.

COROLLARY. --- If \(f\) is convex in \(\mathbf{X}\) , then for all \(\alpha \in \mathbf{R}\) , the set of \(x \in \mathbf{X}\) such that \(f(x) \leqslant \alpha\) (resp. \(f(x) < \alpha\) ) is convex.

In fact, it is the projection on E of the intersection of F (resp. F') and the hyperplane \(E \times \{\alpha\}\) in \(E \times R\) .

PROPOSITION 20. --- Let f be a convex function, defined over a convex set X of the affine space E. Then for every family \((x_{i})_{1\leqslant i\leqslant p}\) of \(p\) points of X and every family \((\lambda_{i})_{1\leqslant i\leqslant p}\) . of \(p\) real numbers, all \(\geqslant 0\) , such that \(\sum_{i = 1}^{p}\lambda_{i} = 1\) , we have:

\[
f \left(\sum_ {i = 1} ^ {p} \lambda_ {i} x _ {i}\right) \leqslant \sum_ {i = 1} ^ {p} \lambda_ {i} f \left(x _ {i}\right).\tag{3}
\]

If \(f\) is strictly convex and if \(\lambda_i > 0\) for all \(i\) , then

\[
f \left(\sum_ {i = 1} ^ {p} \lambda_ {i} x _ {i}\right) <   \sum_ {i = 1} ^ {p} \lambda_ {i} f \left(x _ {i}\right),\tag{4}
\]

unless all the \(x_{i}\) are equal.

The inequality (3) follows from II, prop. 19 above and II, p.~8, prop. 1. Suppose that the \(x_{i}\) are not all equal (which implies \(p \geqslant 2\) ) and that the \(\lambda_{i}\) are all \(>0\) ; then the point \(z = \sum_{i=1}^{p} \lambda_{i} x_{i}\) differs from at least one \(x_{i}\) . Suppose for definiteness that \(z \neq x_{1}\) , write \(z = \lambda_{1} x_{1} + (1 - \lambda_{1}) y_{1}\) where \(y_{1} = \sum_{i=2}^{p} \frac{\lambda_{i}}{1 - \lambda_{1}} x_{i}\) . Then \(y_{1} \neq x_{1}\) and, as \(0 < \lambda_{1} < 1\) , we have, by hypothesis,

\[
f (z) <   \lambda_ {1} f (x _ {1}) + (1 - \lambda_ {1}) f (y _ {1}).
\]

But by (3) \(f(y_{1}) \leqslant \sum_{i=2}^{p} \frac{\lambda_{i}}{1 - \lambda_{1}} f(x_{i})\) , and the inequality (4) follows.

